This bachelor's thesis addresses the automated detection of components in historical handwritten nomenclators using artificial intelligence methods. Nomenclators are cipher systems that combine the substitution of letters, n-grams, nulls, and codes. Manual analysis of such documents is time-consuming and costly, which motivates research into automated approaches.

The aim of this thesis was to design a system for detecting parts of cipher keys using YOLO models, specifically YOLO11. The work included preparing a dataset of 89 digitized pages of historical documents, annotating them into 3 classes representing different nomenclator components, and conducting experiments with various data preprocessing strategies.

The main outcomes of this thesis are trained YOLO11 models and their comparative evaluation. The analysis of the results contributes to ongoing research in nomenclator processing. Detecting individual nomenclator components enables subsequent automated processing, thereby supporting more efficient research in the history of cryptography. At the same time, this work provides a foundation for further improvements in deep-learning methods for processing historical handwritten manuscripts.
